\documentclass[12pt]{article}
\usepackage[utf8]{inputenc}
\usepackage[T1]{fontenc}
\usepackage[margin=1in]{geometry}
\usepackage{setspace}
\usepackage{xcolor}
\usepackage[colorlinks=true, urlcolor=blue!60!black, linkcolor=blue!60!black]{hyperref}
\setlength{\parindent}{0em}
\setlength{\parskip}{0.6em}
\renewcommand{\baselinestretch}{1.0}
\begin{document}
\begin{center}
{\LARGE \textit{Ziglar v. Abbasi}, 582 U.S. 120 (2017)\par}
\vspace{0.5em}
{\large Case Brief --- Two-Page Edition\par}
\end{center}

\section{Memory Jogger}
Post-9/11 detainees could not sue top officials under \textit{Bivens}: extending it is disfavored, and new-context claims fail if special factors counsel hesitation (\href{https://www.law.cornell.edu/supremecourt/text/15-1358}{Cornell LII}).

\section{Detailed Case Facts}
After 9/11, six men of Arab or South Asian descent were held for months in a Brooklyn maximum-security unit under a hold-until-cleared policy, allegedly confined 23 hours a day under constant light, strip-searched, and abused by guards with the wardens' knowledge.

\section{Procedural History}
They sued Attorney General Ashcroft, FBI Director Mueller, INS Commissioner Ziglar, and two wardens. The Second Circuit allowed most claims to proceed.

\section{Judicial Votes}
Reversed in part and remanded in part, 4--2. Kennedy wrote for Roberts, Thomas, and Alito; Thomas concurred in part; Breyer, with Ginsburg, dissented; three Justices did not participate.

\section{Holding}
A claim arises in a new context if it differs meaningfully from \textit{Bivens}, \textit{Davis}, or \textit{Carlson}; then no remedy may be implied if special factors counsel hesitation. The detention-policy claims were barred, the abuse claim against Warden Hasty was remanded, and qualified immunity barred the \S~1985(3) conspiracy claim.

\section{Analysis of Opinions}
\textbf{Majority:} Implied causes of action belong to an ``ancien regime''; Congress usually should decide. High-level policy, national security, Congress's silence after extensive oversight, and alternative remedies all counselled hesitation, and ``even a modest extension is still an extension.'' \textbf{Thomas:} confine \textit{Bivens} to its facts and rethink qualified immunity. \textbf{Breyer:} the claims fit \textit{Bivens} and \textit{Carlson}, Congress ratified \textit{Bivens}, and existing safeguards protect officials. \emph{Later:} \textit{Egbert} (2022) collapsed the test into one question.

\section{Examples: Future Applications}
\textbf{Same-Side:} travelers suing the DHS Secretary over emergency device searches have no remedy. \textbf{Opposite-Side:} a homeowner punched by an ATF agent during a warrantless entry likely does, as in \textit{Bivens}. \textbf{Fence-Sitter:} a pretrial detainee denied insulin by a Marshal.

\section{Critique}
Federal officers now escape damages that state officers face under \S~1983, and the alternative remedies are often illusory; critics say the Court is overruling \textit{Bivens} in slow motion.

\section{Key Quotations}
\begin{quote}
``The proper test for determining whether a case presents a new \textit{Bivens} context is as follows. If the case is different in a meaningful way from previous \textit{Bivens} cases decided by this Court, then the context is new.''
\end{quote}
\noindent --- Justice Kennedy, for the Court, 582 U.S. at 139.

\end{document}
